% G-CEZO POH checklists and emergency procedures, transcribed verbatim (see G-CEZO_POH_Checklists.tex header).
% Included by G-CEZO_POH_Checklists.tex and G-CEZO_Checklist.tex. Do not edit wording without approval (CLAUDE.md).
%=====================================================================
\clearpage
\section{Emergency Procedures (POH Section 3)}
\markboth{Emergency Procedures}{}

\subsection*{3.2 AIRSPEEDS FOR SAFE OPERATIONS}
\subsubsection*{3.2a STALL SPEEDS}
\chk{2325 lbs (0\textdegree\ Flap)}{50 KIAS}
\chk{2325 lbs (Full Flap}{44 KIAS}
\subsubsection*{3.2b MANEUVERING SPEEDS}
\chk{2325 lbs}{111 KIAS}
\chk{1531 lbs}{88 KIAS}
\subsubsection*{3.2c NEVER EXCEED SPEED}
\chk{Never Exceed Speed}{160 KIAS}
\subsubsection*{3.2d POWER OFF GLIDE SPEED}
\chk{2325 lbs (0\textdegree\ Flap)}{73 KIAS}

\subsection*{\color{emerg}3.3 EMERGENCY PROCEDURES CHECKLIST}

\begin{proc}{\color{emerg}3.3a ENGINE FIRE DURING START (3.7)}
\chk{Starter}{CRANK ENGINE}
\chk{Mixture}{IDLE CUT-OFF}
\chk{Throttle}{OPEN}
\chk{Electric Fuel Pump}{OFF}
\chk{Fuel Selector}{OFF}
\txt{Abandon if fire continues}
\end{proc}

\begin{proc}{\color{emerg}3.3b ENGINE POWER LOSS DURING TAKEOFF (3.9)}
\txt{If sufficient runway remains for a normal landing, land straight ahead.}
\cond{If insufficient runway remains:}
\txt{Maintain safe airspeed}
\txt{Make only shallow turn to avoid obstructions}
\txt{Flaps as situation requires}
\cond{If sufficient altitude has been gained to attempt a restart:}
\txt{Maintain safe airspeed}
\chk{Fuel Selector}{SWITCH to tank containing fuel}
\chk{Electric Fuel Pump}{CHECK ON}
\chk{Mixture}{CHECK RICH}
\chk{Carburetor Heat}{ON}
\chk{Primer}{LOCKED}
\txt{If power is not regained, proceed with power off landing (3.3d).}
\end{proc}

\begin{proc}{\color{emerg}3.3c ENGINE POWER LOSS IN FLIGHT (3.11)}
\chk{Fuel Selector}{SWITCH to tank containing fuel}
\chk{Electric Fuel Pump}{ON}
\chk{Mixture}{RICH}
\chk{Carburetor Heat}{ON}
\chk{Engine Gauges}{CHECK for indication of cause of power loss}
\chk{Primer}{CHECK LOCKED}
\txt{If no fuel pressure is indicated, check tank selector position to be sure it is on a tank containing fuel.}
\cond{When power is restored:}
\chk{Carburetor Heat}{OFF}
\chk{Electric Fuel Pump}{OFF}
\cond{If power is not restored, prepare for power off landing (3.3d).}
\txt{Trim for 73 KIAS}
\end{proc}

\begin{proc}{\color{emerg}3.3d POWER OFF LANDING (3.13)}
\txt{Locate suitable field.}
\txt{Establish spiral pattern.}
\txt{1000 ft. above field at downwind position for normal landing approach. When field can easily be reached slow to 63 KIAS for shortest landing.}
\cond{Touchdowns should normally be made at lowest possible airspeed with full flaps extended.}
\cond{When committed to landing:}
\chk{Ignition}{OFF}
\chk{BATT MASTR Switch}{OFF}
\chk{ALTR Switch}{OFF}
\chk{Fuel Selector}{OFF}
\chk{Mixture}{IDLE CUT-OFF}
\chk{Seat Belts and Harnesses}{TIGHT}
\end{proc}

\begin{proc}{\color{emerg}3.3e FIRE IN FLIGHT (3.15)}
\chk{Source of Fire}{CHECK}
\cond{\textbf{Electrical fire (smoke in cabin):}}
\chk{BATT MASTR Switch}{OFF}
\chk{ALTR Switch}{OFF}
\chk{Vents}{OPEN}
\chk{Cabin Heat}{OFF}
\txt{Land as soon as practical.}
\cond{\textbf{Engine fire:}}
\chk{Fuel Selector}{OFF}
\chk{Throttle}{CLOSED}
\chk{Mixture}{IDLE CUT-OFF}
\chk{Electric Fuel Pump}{CHECK OFF}
\chk{Heater}{OFF}
\chk{Defroster}{OFF}
\txt{Proceed with POWER OFF LANDING procedure (3.3d).}
\pohnote{NOTE}{The possibility of an engine fire in flight is extremely remote. The procedure given is general and pilot judgement should be the determining factor for action in such an emergency.}
\end{proc}

\begin{proc}{\color{emerg}3.3f LOSS OF OIL PRESSURE (3.17)}
\txt{Land as soon as possible and investigate cause.}
\txt{Prepare for power off landing (3.3d).}
\end{proc}

\begin{proc}{\color{emerg}3.3g LOSS OF FUEL PRESSURE (3.19)}
\chk{Electric Fuel Pump}{ON}
\chk{Fuel Selector}{CHECK on TANK CONTAINING FUEL}
\end{proc}

\begin{proc}{\color{emerg}3.3h HIGH OIL TEMPERATURE (3.21)}
\txt{Land at nearest airport and investigate the problem.}
\txt{Prepare for power off landing (3.3d).}
\end{proc}

\begin{proc}{\color{emerg}3.3i ELECTRICAL FAILURES (3.23)}
\pohnote{NOTE}{When operating with light electrical load and a fully charged battery, the Alternator Inop. light may illuminate due to minimal alternator output. If the alternator is functional a slight increase in electrical load should extinguish the Inop. indication.}
\cond{ALT annunciator light illuminated:}
\chk{Ammeter}{CHECK to VERIFY inop. alt.}
\cond{If ammeter shows zero:}
\chk{ALTR switch}{OFF}
\cond{Reduce electrical loads to minimum:}
\chk{ALTNTR. FIELD Circuit Breaker}{CHECK and RESET as required}
\chk{ALTR Switch}{ON}
\cond{If power not restored:}
\chk{ALTR Switch}{OFF}
\txt{If alternator output cannot be restored, reduce electrical loads and land as soon as practical. The battery is the only remaining source of electrical power.}
\end{proc}

\begin{proc}{\color{emerg}3.3j ELECTRICAL OVERLOAD (Alternator over 20 amps above known electrical load) (3.25)}
\chk{ALTR Switch}{ON}
\chk{BATT MASTR Switch}{OFF}
\cond{If alternator loads are reduced:}
\chk{Electrical Load}{REDUCE to Minimum}
\txt{Land as soon as practical.}
\pohnote{NOTE}{Due to increased system voltage and radio frequency noise, operation with ALT switch ON and BATT switch OFF should be made only when required by an electrical system failure.}
\cond{If alternator loads are not reduced:}
\chk{ALTR switch}{OFF}
\chk{BATT MASTR Switch}{AS REQUIRED}
\txt{Land as soon as possible. Anticipate complete electrical failure.}
\end{proc}

\begin{proc}{\color{emerg}3.3k SPIN RECOVERY (3.27)}
\chk{Throttle}{IDLE}
\chk{Ailerons}{NEUTRAL}
\chk{Rudder}{FULL OPPOSITE to DIRECTION of ROTATION}
\chk{Control Wheel}{FULL FORWARD}
\chk{Rudder}{NEUTRAL (when rotation stops)}
\chk{Control Wheel}{AS REQUIRED to SMOOTHLY REGAIN LEVEL FLIGHT ATTITUDE}
\end{proc}

\begin{proc}{\color{emerg}3.3m OPEN DOOR (3.29)}
\txt{If both upper and lower latches are open, the door will trail slightly open and airspeeds will be reduced slightly.}
\cond{To close the door in flight:}
\txt{Slow airplane to 89 KIAS}
\chk{Cabin Vents}{CLOSE}
\chk{Storm Window}{OPEN}
\chk{If upper latch is open}{LATCH}
\chk{If side latch if open}{PULL on ARMREST while moving latch handle to LATCH position.}
\chk{If both latches are open}{LATCH SIDE LATCH then TOP LATCH.}
\end{proc}

\begin{proc}{\color{emerg}3.3n ENGINE ROUGHNESS (3.31)}
\chk{Carburetor Heat}{ON}
\cond{If roughness continues after one minute:}
\chk{Carburetor Heat}{OFF}
\chk{Mixture}{ADJUST for MAXIMUM SMOOTHNESS}
\chk{Electric Fuel Pump}{ON}
\chk{Fuel Selector}{SWITCH TANKS}
\chk{Engine Gauges}{CHECK}
\chk{Magneto Switch}{L then R then BOTH}
\cond{If operation is satisfactory on either magneto, continue on that magneto at reduced power and full RICH mixture to first airport.}
\cond{Prepare for power off landing (3.3d).}
\pohnote{NOTE}{Partial carburetor heat may be worse than no heat at all, since it may melt part of the ice which will refreeze in the intake system. Therefore, when using carburetor heat always use full heat; and, when ice is removed, return the control to the full cold position.}
\end{proc}

\begin{proc}{\color{emerg}3.3o CARBURETOR ICING (3.33)}
\chk{Carburetor Heat}{ON}
\chk{Mixture}{ADJUST for MAXIMUM SMOOTHNESS}
\end{proc}

%=====================================================================
\clearpage
\section{Normal Procedures (POH Section 4)}
\markboth{Normal Procedures}{}

\subsection*{4.3 AIRSPEEDS FOR SAFE OPERATIONS}
\chk{(a) Best Rate of Climb Speed}{79 KIAS}
\chk{(b) Best Angle of Climb Speed}{63 KIAS}
\chk{(c) Turbulent Air Operating Speed (See Subsection 2.3)}{111 KIAS}
\chk{(d) Maximum Flap Speed}{103 KIAS}
\chk{(e) Landing Final Approach Speed (Flaps 40\textdegree)}{63 KIAS}
\chk{(f) Maximum Demonstrated Crosswind Velocity}{17 KTS}

\subsection*{\color{normal}4.5 NORMAL PROCEDURES CHECKLIST}

\begin{proc}{\color{normal}4.5a Preparation (4.9)}
\chk{Airplane Status}{AIRWORTHY, PAPERS ON BOARD}
\chk{Weather}{SUITABLE}
\chk{Baggage}{WEIGHED, STOWED, TIED}
\chk{Weight and C.G.}{WITHIN LIMITS}
\chk{Navigation}{PLANNED}
\chk{Charts and Navigation Equipment}{ON BOARD}
\chk{Performance and Range}{COMPUTED AND SAFE}
\pohnote{\textit{CAUTION}}{The flap position should be noted before boarding the airplane. The flaps must be placed in the UP position before they will lock and support weight on the step.}
\end{proc}

\begin{proc}{\color{normal}4.5b Preflight Checklists (4.11)}
\subsubsection*{COCKPIT (4.11a)}
\chk{Control Wheel}{RELEASE BELTS}
\chk{Circuit Breakers}{IN}
\chk{Avionics}{OFF}
\chk{Parking Brake}{SET}
\chk{Electric Switches}{OFF}
\chk{Magneto Switch}{OFF}
\chk{Mixture}{IDLE CUT-OFF}
\chk{BATT MASTR Switch}{ON}
\chk{Fuel Quantity Gauges}{CHECK}
\chk{Annunciator Panel}{CHECK}
\chk{BATT MASTR Switch}{OFF}
\chk{Flight Controls}{CHECK}
\chk{Flaps}{CHECK}
\chk{Trim}{CHECK, SET NEUTRAL}
\chk{Pitot Drain}{DRAIN, CLOSE}
\chk{Static Drain}{DRAIN, CLOSE}
\chk{Windows}{CHECK, CLEAN}
\chk{Tow Bar}{STOW}
\chk{Baggage}{SECURE}

\subsubsection*{RIGHT WING (4.11b)}
\chk{Wing}{FREE of ICE, SNOW, FROST}
\chk{Control Surfaces}{CHECK for interference - CLEAR of ICE, SNOW, FROST}
\chk{Hinges}{CHECK for interference}
\chk{Static Wicks}{CHECK}
\chk{Wing Tip and Lights}{CHECK}
\chk{Fuel Tank}{CHECK supply visually - SECURE CAPS}
\pohnote{\textit{CAUTION}}{When draining any amount of fuel, care should be taken to ensure that no fire hazard exists before starting engine.}
\chk{Fuel Tank Sump}{DRAIN, CHECK for water, sediment and proper fuel}
\chk{Fuel Vent}{CLEAR}
\chk{Tie Down and Chock}{REMOVE}
\chk{Main Gear Strut}{PROPER INFLATION (4.50 in.)}
\chk{Tire}{CHECK}
\chk{Brake Block and Disc}{CHECK}
\chk{Fresh Air Inlet}{CLEAR}

\subsubsection*{NOSE SECTION (4.11c)}
\chk{Fuel and Oil}{CHECK FOR LEAKS}
\chk{Cowling}{SECURE}
\chk{Windshield}{CLEAN}
\chk{Propeller and Spinner}{CHECK}
\chk{Air Inlets}{CLEAR}
\chk{Alternator Belt}{CHECK TENSION}
\chk{Landing Light}{CHECK}
\chk{Nose Chock}{REMOVE}
\chk{Nose Gear Strut}{PROPER INFLATION (3.25 in.)}
\chk{Nose Wheel Tire}{CHECK}
\chk{Oil}{CHECK LEVEL}
\chk{Dipstick}{PROPERLY SEATED}
\chk{Fuel Strainer}{DRAIN, CHECK for water, sediment and proper fuel}

\subsubsection*{LEFT WING (4.11d)}
\chk{Wing}{FREE of ICE, SNOW, FROST}
\chk{Fresh Air Inlet}{CLEAR}
\chk{Main Gear Strut}{PROPER INFLATION (4.50 in.)}
\chk{Tire}{CHECK}
\chk{Brake Block and Disc}{CHECK}
\chk{Fuel Tank Sumps}{DRAIN, CHECK for water, sediment and proper fuel}
\chk{Fuel Vent}{CLEAR}
\chk{Tie Down and Chock}{REMOVE}
\chk{Fuel Tanks}{CHECK supply visually - SECURE CAPS}
\chk{Pitot/Static Head}{REMOVE COVER - HOLES CLEAR}
\chk{Wing Tip and Lights}{CHECK}
\chk{Control Surfaces}{CHECK for interference - FREE of ICE, SNOW, FROST}
\chk{Hinges}{CHECK for interference}
\chk{Static Wicks}{CHECK}

\subsubsection*{FUSELAGE (4.11e)}
\chk{Antennas}{CHECK}
\chk{Empennage}{FREE of ICE, SNOW, FROST}
\chk{Fresh Air Inlet}{CLEAR}
\chk{Stabilator and Trim Tab}{CHECK for interference}
\chk{Tie Down}{REMOVE}
\chk{BATT MASTR Switch}{ON}
\chk{Cockpit Lighting}{CHECK}
\chk{Nav and Strobe Lights}{CHECK}
\chk{Stall Warning}{CHECK}
\chk{Pitot Heat}{CHECK}
\chk{All Switches}{OFF}
\chk{Passengers}{BOARD}
\chk{Cabin Door}{CLOSE and SECURE}
\chk{Seat Belts and Harnesses}{FASTEN - CHECK interia reel}
\end{proc}

\begin{proc}{\color{normal}4.5c Before Starting Engine Checklist (4.13)}
\subsubsection*{BEFORE STARTING ENGINE (4.13)}
\chk{Brakes}{SET}
\chk{Carburetor Heat}{FULL OFF}
\chk{Fuel Selector}{DESIRED TANK}
\chk{Avionics}{OFF}
\chk{ALTR Switch}{ON}
\end{proc}

\begin{proc}{\color{normal}4.5d Engine Start Checklist (4.15)}
\subsubsection*{NORMAL START - COLD ENGINE (4.15a)}
\chk{Throttle}{\textonequarter\ INCH OPEN}
\chk{BATT MASTR Switch}{ON}
\chk{Electric Fuel Pump}{ON}
\chk{Mixture}{FULL RICH}
\chk{Starter}{ENGAGE}
\chk{Throttle}{ADJUST}
\chk{Oil Pressure}{CHECK}
\txt{If engine does not start within 10 sec., prime and repeat starting procedure.}

\subsubsection*{NORMAL START - HOT ENGINE (4.15b)}
\chk{Throttle}{\textonehalf\ INCH OPEN}
\chk{BATT MASTR Switch}{ON}
\chk{Electric Fuel Pump}{ON}
\chk{Mixture}{FULL RICH}
\chk{Starter}{ENGAGE}
\chk{Throttle}{ADJUST}
\chk{Oil Pressure}{CHECK}

\subsubsection*{ENGINE START WHEN FLOODED (4.15c)}
\chk{Throttle}{OPEN FULL}
\chk{BATT MASTR Switch}{ON}
\chk{Electric Fuel Pump}{OFF}
\chk{Mixture}{IDLE CUT-OFF}
\chk{Starter}{ENGAGE}
\chk{Mixture}{ADVANCE}
\chk{Throttle}{RETARD}
\chk{Oil Pressure}{CHECK}

\subsubsection*{ENGINE START WITH EXTERNAL POWER SOURCE (4.15d)}
\chk{BATT MASTR Switch}{OFF}
\chk{All Electrical Equipment}{OFF}
\chk{Terminals}{CONNECT}
\chk{External Power Plug}{INSERT in receptacle}
\pohnote{NOTE}{For all normal operations using the PEP jumper cables, the battery master switch should be OFF. However, it is possible to use the ship's battery in parallel by turning the battery master switch ON. This will give longer cranking capabilities, but will not increase the amperage.}
\pohnote{\textit{CAUTION}}{Care should be exercised, because, if the ship's battery has been depleted, the external power supply can be reduced to the level of the ship's battery. This can be tested by turning the battery master switch ON momentarily while the starter is engaged. If cranking speed increases, the ship's battery is at a higher level than the external power supply.}
\txt{Proceed with normal start}
\chk{Throttle}{LOWEST POSSIBLE RPM}
\chk{External Power Plug}{DISCONNECT from receptacle}
\chk{BATT MASTR Switch}{ON - CHECK AMMETER}
\chk{Oil Pressure}{CHECK}
\end{proc}

\begin{proc}{\color{normal}4.5e Engine Warm-Up Checklist (4.17)}
\subsubsection*{ENGINE WARM-UP (4.17)}
\chk{Throttle}{800 to 1200 RPM}
\end{proc}

\begin{proc}{\color{normal}4.5f Taxiing Checklist (4.19)}
\subsubsection*{TAXIING (4.19)}
\chk{Parking Brake}{RELEASE}
\chk{Taxi Area}{CLEAR}
\chk{Throttle}{APPLY SLOWLY}
\chk{Brakes}{CHECK}
\chk{Steering}{CHECK}
\end{proc}

\begin{proc}{\color{normal}4.5g Ground Check Checklist (4.21)}
\subsubsection*{GROUND CHECK (4.21)}
\chk{Parking Brake}{SET}
\chk{Throttle}{2000 RPM}
\chk{Magnetos}{max. drop 175 RPM -max. diff. 50 RPM}
\chk{Vacuum}{4.8 to 5.1 in. Hg}
\chk{Oil Temperature}{CHECK}
\chk{Oil Pressure}{CHECK}
\chk{Air Conditioner}{CHECK}
\chk{Annunciator Panel}{PRESS-TO-TEST}
\chk{Carburetor Heat}{CHECK}
\txt{Engine is warm for takeoff when throttle can be opened without engine faltering.}
\chk{Electric Fuel Pump}{OFF}
\chk{Fuel Pressure}{CHECK}
\chk{Throttle}{RETARD}
\end{proc}

\begin{proc}{\color{normal}4.5h Before Takeoff Checklist (4.23)}
\subsubsection*{BEFORE TAKEOFF (4.23)}
\chk{Battery Master Switch}{ON}
\chk{Alternator Switch}{ON}
\chk{Flight Instruments}{CHECK}
\chk{Fuel Selector}{PROPER TANK}
\chk{Electric Fuel Pump}{ON}
\chk{Engine Gauges}{CHECK}
\chk{Carburetor Heat}{OFF}
\chk{Seat Backs}{ERECT}
\pohnote{NOTE}{The mixture should be set FULL RICH, but a minimum amount of leaning is permitted for smooth engine operation when taking off at high elevation.}
\chk{Mixture}{SET}
\chk{Primer}{LOCK}
\pohnote{NOTE}{The inertia reel type shoulder harness should be given a pull test to check its locking restraint feature.}
\pohnote{NOTE}{If the fixed shoulder harness (non-inertial reel type) are installed for the aft seat occupants they should be adjusted to provide adequate restraint.}
\chk{Belts Harness}{FASTENED CHECK}
\chk{Empty Seats}{SEAT BELTS SNUGLY FASTENED}
\chk{Flaps}{SET}
\chk{Trim}{SET}
\chk{Controls}{FREE}
\chk{Door}{LATCHED}
\chk{Air Conditioner}{OFF}
\end{proc}

\begin{proc}{\color{normal}4.5i Takeoff Checklist (4.25)}
\subsubsection*{NORMAL (4.25a)}
\chk{Flaps}{SET}
\chk{Trim}{SET}
\txt{Accelerate to 45 to 55 KIAS}
\chk{Control Wheel}{back pressure to rotate to climb attitude}

\subsubsection*{0\textdegree\ FLAPS TAKEOFF PERFORMANCE (4.25b)}
\chk{Flaps}{UP}
\txt{Accelerate to 40-50 KIAS (depending on weight)}
\chk{Control Wheel}{back pressure to rotate to climb attitude}
\cond{Accelerate to and maintain 44 to 55 KIAS (depending on weight) until obstacle clearance is achieved and climb out at 79 KIAS.}

\subsubsection*{25\textdegree\ FLAPS TAKEOFF PERFORMANCE (4.25c)}
\chk{Flaps}{25\textdegree\ (second notch)}
\txt{Accelerate to 40-50 KIAS (depending on weight)}
\chk{Control Wheel}{back pressure to rotate to climb attitude}
\cond{Accelerate to and maintain 44 to 55 KIAS (depending on weight) until obstacle clearance is achieved and climb out at 79 KIAS.}
\chk{Flaps}{RETRACT SLOWLY}

\subsubsection*{SOFT FIELD, OBSTACLE CLEARANCE (4.25d)}
\chk{Flaps}{25\textdegree\ (second notch)}
\cond{Accelerate and lift off nose gear as soon as possible. Lift off at lowest possible airspeed. Accelerate just above ground to 52 KIAS to climb past obstacle height. Continue climbing while accelerating to best rate of climb speed, 79 KIAS.}
\chk{Flaps}{RETRACT SLOWLY}

\subsubsection*{SOFT FIELD, NO OBSTACLE (4.25e)}
\chk{Flaps}{25\textdegree\ (second notch)}
\cond{Accelerate and lift off nose gear as soon as possible. Lift off at lowest possible airspeed. Accelerate just above ground to best rate of climb speed, 79 KIAS.}
\chk{Flaps}{RETRACT SLOWLY}
\end{proc}

\begin{proc}{\color{normal}4.5j Climb Checklist (4.27)}
\subsubsection*{CLIMB (4.27)}
\chk{Best Rate Climb Speed (Flaps Up)}{79 KIAS}
\chk{Best Angle Climb Speed (Flaps Up)}{63 KIAS}
\chk{En Route}{87 KIAS}
\chk{Electric Fuel Pump}{OFF at desired altitude}
\end{proc}

\begin{proc}{\color{normal}4.5k Cruise Checklist (4.29)}
\subsubsection*{CRUISE}
\txt{Reference performance charts and Avco-Lycoming Operators Manual.}
\chk{Normal Maximum Power}{75\%}
\chk{Power}{SET per power table}
\chk{Mixture}{ADJUST}
\end{proc}

\begin{proc}{\color{normal}4.5m Descent Checklist (4.31)}
\subsubsection*{NORMAL DESCENT (4.31a)}
\chk{Throttle}{2500 rpm}
\chk{Airspeed}{126 KIAS}
\chk{Mixture}{RICH}
\chk{Carburetor Heat}{ON if required}
\subsubsection*{POWER OFF DESCENT (4.31b)}
\chk{Carburetor Heat}{ON if required}
\chk{Throttle}{CLOSE}
\chk{Airspeed}{AS REQUIRED}
\chk{Mixture}{AS REQUIRED}
\chk{Power}{VERIFY with throttle EVERY 30 SECONDS}
\end{proc}

\begin{proc}{\color{normal}4.5n Approach And Landing Checklist (4.33)}
\subsubsection*{APPROACH AND LANDING (See charts in Section 5) (4.33)}
\chk{Fuel Selector}{PROPER TANK}
\chk{Seat Backs}{ERECT}
\chk{Belts/Harness}{FASTEN/CHECK}
\pohnote{NOTE}{If the fixed shoulder harness (non-inertia reel type) is installed, it must be connected to the seat belt and adjusted to allow proper accessibility to all controls including fuel selector, flaps, trim, etc, while maintaining adequate restraint for the occupant.}
\pohnote{NOTE}{If the inertia reel type shoulder harness is installed, a pull test of its locking restraint feature should be performed.}
\chk{Electric Fuel Pump}{ON}
\chk{Mixture}{SET}
\chk{Flaps}{SET - 103 KIAS MAXIMUM}
\chk{Air Conditioner}{OFF}
\txt{Trim to 70 KIAS}
\chk{Final Approach Speed (Flaps 40\textdegree)}{63 KIAS}
\end{proc}

\begin{proc}{\color{normal}4.5o Stopping Engine Checklist (4.35)}
\subsubsection*{STOPPING ENGINE (4.35)}
\chk{Flaps}{RETRACT}
\chk{Electric Fuel Pump}{OFF}
\chk{Electrical Switches}{OFF}
\chk{Air Conditioner}{OFF}
\chk{Radios}{OFF}
\chk{Throttle}{FULL AFT}
\chk{Mixture}{IDLE CUT-OFF}
\chk{Magnetos}{OFF}
\chk{ALTR Switch}{OFF}
\chk{BATT MASTR Switch}{OFF}
\chk{Trim}{SET}
\pohnote{NOTE}{When alternate fuels are used, the engine should be run up to 1200 rpm for one minute prior to shutdown to clean out any unburned fuel.}
\pohnote{NOTE}{The flaps must be placed in the UP position for the flap step to support weight. Passengers should be cautioned accordingly.}
\end{proc}

\begin{proc}{\color{normal}4.5p Mooring Checklist (4.37)}
\subsubsection*{MOORING (4.37)}
\chk{Parking Brake}{SET}
\chk{Control Wheel}{SECURE with belts}
\chk{Flaps}{FULL UP}
\chk{Wheel Chocks}{IN PLACE}
\chk{Tie Downs}{SECURE}
\end{proc}

%=====================================================================
\clearpage
\section{Cockpit Placards (POH 2.23)}
\markboth{Placards}{}

\begin{proc}{TAKEOFF CHECKLIST}
\txt{In full view of the pilot:}
\begin{multicols}{2}
\txt{Fuel on proper tank}
\txt{Electric fuel pump on}
\txt{Engine gauges checked}
\txt{Flaps - set}
\txt{Carb. heat off}
\txt{Mixture set}
\txt{Primer locked}
\txt{Seat backs erect}
\txt{Fasten belts/harness}
\txt{Trim tab - set}
\txt{Controls - free}
\txt{Door - latched}
\txt{Air conditioner off}
\end{multicols}
\end{proc}

\begin{proc}{LANDING CHECKLIST}
\begin{multicols}{2}
\txt{Fuel on proper tank}
\txt{Mixture rich}
\txt{Electric fuel pump on}
\txt{Seat backs erect}
\txt{Flaps - set (White Arc)}
\txt{Fasten belts/harness}
\txt{Air conditioner off}
\end{multicols}
\txt{In air conditioned aircraft, the AIR COND OFF item in the takeoff and landing checklists above, is mandatory.}
\end{proc}

